\documentclass[12pt]{article}
\usepackage[utf8]{inputenc}
\usepackage{amsmath,amssymb,latexsym}
\usepackage{enumerate,comment}
\usepackage{graphicx}
\usepackage[spanish]{babel}
\usepackage{xypic}

\textheight=25cm \textwidth=18cm \topmargin=-2cm \oddsidemargin=-1cm

\newcommand{\abs}[1]{\left\vert#1\right\vert}
\newcommand{\norm}[1]{\left\Vert#1\right\Vert}

\newcommand{\Co}{\mathfrak C}
\newcommand{\diag}{\mathrm{diag}}
\newcommand{\Ext}{\mathrm{ext}}
\newcommand{\F}{\mathfrak F}
\newcommand{\h}{\mathrm h}
\newcommand{\Int}{\mathrm{int}}
\newcommand{\n}{\mathfrak{n}}
\newcommand{\N}{\mathbb N}
\newcommand{\Q}{\mathbb Q}
\newcommand{\R}{\mathbb R}
\newcommand{\US}{\mathbb{S}^2}
\newcommand{\x}{\mathrm x}
\newcommand{\xk}{(\x_k)_{k\in \N}}
\newcommand{\y}{\mathrm y}
\newcommand{\z}{\mathrm z}
\newcommand{\Z}{\mathbb Z}

\newcommand{\eps}{\varepsilon}
\newcommand{\lbil}{\left <}
\newcommand{\rbil}{\right >}
\newcommand{\To}{\longrightarrow}
\newcommand{\mTo}{\longmapsto}

\begin{document}
\pagestyle{empty}
\begin{center}
{\textsc{Cálculo Diferencial e Integral III} \\
\medskip 
\large{Ejercicios I - 1 \\
(El espacio euclidiano de dimensión $n$)}}
\end{center}
\bigskip

En general, $\Omega,\Delta, \dots$ serán subconjuntos del espacio euclidiano $\R^n$.

\begin{enumerate}

\item 
\begin{enumerate}[(1)]
\item Verifica que las funciones
\[ 
\norm{\x}_{\max} := \max\{\abs{x_1},\ldots,\abs{x_n}\} \quad \text{y}\quad 
\norm{\x}_1 := \abs{x_1} + \cdots + \abs{x_n} 
\]
satisfacen las propiedades de una norma en $\R^n$.

\textbf{Solución: } Veamos que \(\norm{\x}_{\max}\) es una norma:
    \begin{itemize}
        \item Por definición del valor absoluto, \(\abs{x_i} \geq 0\) para toda \(i \in [n]\), y ya que \(\max\{\abs{x_1}, \dots, \abs{x_n}\} \geq \abs{x_i}\), se sigue que \(\max\{\abs{x_1}, \dots, \abs{x_n}\} \geq 0\). Por lo tanto, \(\norm{\x}_{\max} \geq 0\) para toda \(\x\).
        \item Sea \(\lambda \in \R\). Por definición del producto escalar en \(\R^n\), \(\lambda \x = (\lambda x_1, \dots, \lambda x_n)\). Al aplicar el valor absoluto, obtenemos \(\abs{\lambda x_i} = \abs{\lambda} \abs{x_i}\). Supongamos primero que \(\lambda = 0\), entonces \(\max_{1 \leq i \leq n} \{\lambda x_i\} = \lambda = 0\). Ahora supongamos que \(\lambda \neq 0\), entonces \(\abs{\lambda} > 0\). Ya que multiplicar por un número positivo preserva las desigualdades, \(\max_{1 <= i <= n}\{\abs{\lambda x_i}\} = \abs{\lambda} \max_{1 <= i <= n}\{\abs{x_i}\} = \abs{\lambda} \norm{\x}_{\max}\).
        \item Por definición, \(\norm{\x}_{\max} + \norm{\y}_{\max} = \max_{i\in[n]}\{x_i\} + \max_{i\in[n]}\{y_i\}\). Ya que \(\max_{i\in[n]}\{x_i\} \geq x_j\) y \(\max_{i\in[n]}\{y_i\} \geq y_j\) para cualquier \(j \leq n\), entonces \(\max_{i\in[n]}\{x_i\} + \max_{i\in[n]}\{y_i\} \geq \abs{x_j} + \abs{y_j} \geq \abs{x_j + y_j}\) para cualquier \(j\), lo que implica que \(\max_{i\in[n]}\{x_i\} + \max_{i\in[n]}\{y_i\} \geq \max_{i \in [n]}\{\abs{x_j + y_j}\}\) y por lo tanto, \(\norm{\x} + \norm{\y} \geq \norm{\x + \y}\).
    \end{itemize}
Ahora, veamos que \(\norm{\x}_1\) es una norma:
\begin{itemize}
    \item Por definición del valor absoluto, \(\abs{x_i} \geq 0\) y ya que la suma de valores no negativos es no negativa, \(\norm{\x}_1\) es no negativa para cualquier vector \(\x \in \R^n\).
    
    Además, ya que \(\sum_{i \in [n]}|x_i| = 0\) si y solo si \(x_i = 0\) para cualquier índice \(i\), se sigue que \(\norm{\x}_1 = 0\) si y solo si \(\x = 0\).
    \item Sea \(\lambda \in \mathbb{R}\). Por definición del producto escalar, tenemos que 
    \[\norm{\lambda \x}_1 = \sum_{i = 1}^{n}\abs{\lambda x_i} = \sum_{i = 1}^{n}\abs{\lambda}\abs{x_i} = \abs{\lambda}\sum_{i = 1}^{n}\abs{x_i} = \abs{\lambda}\norm{\x}_1.\]
    \item Reescribimos
    \[
    \norm{\x}_1 + \norm{\y}_1 = \sum_{i = 1}^{n}\abs{x_i} + \sum_{i = 1}^{n}\abs{y_i}.
    \]
    Sumando término a término \(\abs{x_i} + \abs{y_i}\), obtenemos la expresión
    \[
    \sum_{i = 1}^{n}\abs{x_i} + \sum_{i = 1}^{n}\abs{y_i} = \sum_{i = 1}^{n}(\abs{x_i} + \abs{y_i}) \geq \sum_{i = 1}^{n}\abs{x_i + y_i}.
    \]
    Por lo tanto, 
    \[
    \norm{\x}_1 + \norm{\y}_1 \geq \norm{\x + \y}_1.
    \]
    
\end{itemize}
\item Describe las bolas unitarias en $\R^2$ con respecto a las normas $\norm{\x}_{\max}$ y $\norm{\x}_1$. 
\begin{itemize}
    \item La bola unitaria es el conjunto \(B_{\max} = \{(x_1, x_2) \in \R^2 : \norm{\x}_{\max} \leq 1\}\), lo que significa que
    \[
    \max\{\abs{x_1}, \abs{x_2}\} \leq 1
    \]
    que es equivalente a
    \[
    -1 \leq x_1 \leq 1 \quad \text{y} \quad -1 \leq x_2 \leq 1,
    \]
    por lo que la bola unitaria es un cuadrado de lado 2 centrado alrededor del origen:
    \[
    B_{\max} = [-1, 1] \times [-1, 1].
    \]
    \item Para la norma \(\norm{\x}_1\), la bola unitaria se define como el conjunto \(B_1 = \{(x_1, x_2) \in \R^2 : \norm{\x}_1 \leq 1\}\). Es decir, 
    \[
    \abs{x_1} + \abs{x_2} \leq 1.
    \]
    Esto describe un rombo (o un cuadrado rotado) en el origen con vértices que ocurren cuando una coordenada es 0 y la otra es 1: \((1, 0)\), \((0, 1)\), \((0, -1)\) y \((-1, 0)\). 
\end{itemize}
\end{enumerate}

\item Demuestra que se cumplen las siguientes desigualdades:
\begin{enumerate}[(1)]
\item $\norm{\x}_{\max}\leq \norm{\x}\leq \sqrt{n}\norm{\x}_{\max}$.
\item $\frac{1}{\sqrt{n}}\norm{\x}_1\leq \norm{\x}\leq \norm{\x}_1$.
\end{enumerate}

\item Describe el interior de
\[ \Omega = \{ (x,y,z)\in \R^3 : 0\leq x < 1, y^2+z^2\leq 1\}. \]

\item Encuentra la frontera de 
\[ \Omega = \{ (x,y)\in \R^2 : x^2 - y^2 > 1 \}. \]

\item Decir si las siguientes propiedades son ciertas. Prueba tu respuesta.
\begin{enumerate}[(1)]
\item $\Int(\Omega\cup \Delta) = \Int(\Omega)\cup \Int(\Delta),\; \Int(\Omega\cap \Delta) = \Int(\Omega)\cap \Int(\Delta)$.
\item $\Ext(\Omega\cup \Delta) = \Ext(\Omega)\cup \Ext(\Delta),\; \Ext(\Omega\cap \Delta) = \Ext(\Omega)\cap \Ext(\Delta)$.
\end{enumerate}

\item Decir si las siguientes propiedades son ciertas. Prueba tu respuesta.
\begin{enumerate}[(1)]
\item $\partial(\Omega\cup \Delta) = \partial \Omega\cup \partial \Delta,\; 
\partial(\Omega\cap \Delta) = \partial \Omega\cap \partial \Delta$.
\item Si $\Omega\subset \Delta$ entonces $\Int(\Omega)\subset \Int(\Delta),\; \Ext(\Omega)\subset \Ext(\Delta),\; \partial \Omega\subset \partial \Delta$.
\end{enumerate}

\item Define la \emph{suma} de $\Omega$ y $\Delta$ como
\[ \Omega + \Delta := \{ \x+\y\in \R^n : \x\in \Omega, \y\in \Delta \}. \]
Demuestra que $\Omega + \Delta$ es abierto.

\item Encuentra la cerradura de 
\[ \Omega = \{ (x,y)\in \R^2 : x > y^2 \}. \]

\item Prueba que si $\Omega$ es cerrado entonces 
$\Int(\partial \Omega)=\emptyset$.

\item Supongamos que $a_1,\ldots,a_n,b_1,\ldots,b_n\in \R^n$ son tales que $a_j<b_j$ para toda $j=1,\ldots,n$. Hemos definimos el rectángulo cerrado (abierto) por
\[ R = [a_1,b_1]\times \cdots\times [a_n,b_n], \]
o bien,
\[ R = (a_1,b_1)\times \cdots\times (a_n,b_n). \]
En cualquier caso, la \textsf{diagonal} de $R$ se define por:
\[ \diag(R) := \norm{(b_1,\ldots,b_n) - (a_1,\ldots,a_n)}. \]
Prueba que todo rectángulo cerrado $R$ es un subconjunto cerrado, acotado y convexo de $\R^n$. 

\end{enumerate}

\end{document}